The resolution requires two departures from the classical framework, each empirically motivated and individually insufficient. First, not all sellers are informed. At listing time, only 51.8\% of apartment sellers and 41.7\% of house sellers hold a certified energy audit. When buyers observe a missing EPC score, they face genuine ambiguity. The non-disclosing seller may be concealing an unfavorable rating or may simply not know her property's energy performance. Uninformed sellers therefore provide strategic hiders with involuntary cover, preventing silence from being read as the worst possible signal. Second, not all buyers price energy performance into their offers. When the share of energy-sensitive buyers is moderate, the revenue penalty from concealment is diluted, making hiding profitable even against a skeptical pool. Neither departure suffices alone: uninformed sellers keep silence ambiguous, but buyer heterogeneity makes exploiting this ambiguity financially worthwhile. In apartment markets a third ingredient is needed, because no belief about silence can price a silent listing above a verified efficient one: buyers attend to energy much less when a listing shows no rating than when it shows one. % [AUTHOR REWRITE]
